\documentclass[11pt]{article}
%%%%%%%%%%%% PREAMBULO %%%%%%%%%%%%%%%%%%
\usepackage[T1]{fontenc}%indica que usamos la ñ
\usepackage[utf8]{inputenc}%indica el tipo de codificación ISO-8859-1 (latin)ó utf8
\usepackage[spanish]{babel}%indica que escribiremos en español
\title{plantilla para reportes mecatrónica-mecánica UTA}
\usepackage{amsmath}
\usepackage{amssymb,amsfonts,latexsym}
\usepackage{graphicx}
\usepackage{epstopdf}%para convertir figuras a formato pdf
\usepackage{color}
\usepackage{url}  \urlstyle{same}%para conectar url strings en bibliografías
\usepackage{subfigure}%para colocar varias figuras
\usepackage{multicol}
\usepackage{changepage}
\usepackage{float}%para colocar las tablas como H
\usepackage{array}%dar formato a las tablas
\usepackage{longtable}%dar formato a las tablas
\usepackage{bm}
\usepackage{capt-of}
\usepackage{sidecap}
	\sidecaptionvpos{figure}{c}
\usepackage{caption}
\usepackage{commath}
\usepackage{cancel}

%%%%%%%%%%% CONFIGURACION DEL DOCUMENTO %%%%%%%%%%%%%
\usepackage{anysize}%para personalizar el ancho de los margenes
\marginsize{2cm}{2cm}{2cm}{2cm}%izquierda, derecha, arriba, abajo
\usepackage{anyfontsize}

\usepackage{appendix}
\renewcommand{\appendixname}{Apéndices}
\renewcommand{\appendixtocname}{Apéndices}
\renewcommand{\appendixpagename}{Apéndices} 

% Para que las referencias sean hipervínculos a las figuras o ecuaciones y
% aparezcan en color
\usepackage[colorlinks=true,plainpages=true,citecolor=blue,linkcolor=blue]{hyperref}
%\usepackage{hyperref} 
% Para agregar encabezado y pie de página
\usepackage{fancyhdr} 
\pagestyle{fancy}
\fancyhf{}
\fancyhead[L]{\footnotesize UTA} %encabezado izquierda
\fancyhead[R]{\footnotesize Esc. de Ingeniería}   % dereecha
\fancyfoot[R]{\footnotesize Reporte}  % Pie derecha
\fancyfoot[C]{\thepage}  % centro
\fancyfoot[L]{\footnotesize Ing. Mecanica y Mecatrónica}  %izquierda
\renewcommand{\footrulewidth}{0.4pt}

\usepackage{listings} % Para usar código fuente
\definecolor{dkgreen}{rgb}{0,0.6,0} % Definimos colores para usar en el código
\definecolor{gray}{rgb}{0.5,0.5,0.5} 
% configuración para el lenguaje que queramos utilizar
\lstset{language=Matlab,
   keywords={break,case,catch,continue,else,elseif,end,for,function,
      global,if,otherwise,persistent,return,switch,try,while},
   basicstyle=\ttfamily,
   keywordstyle=\color{blue},
   commentstyle=\color{red},
   stringstyle=\color{dkgreen},
   numbers=left,
   numberstyle=\tiny\color{gray},
   stepnumber=1,
   numbersep=10pt,
   backgroundcolor=\color{white},
   tabsize=4,
   showspaces=false,
   showstringspaces=false}

\newcommand{\sen}{\operatorname{\sen}}	% Definimos el comando \sen para el seno
%en español

\title{Plantilla para informes ing. mecanica-mecatronica}

%%%%%%%%%%% COMIENZO DEL DOCUMENTO %%%%%%%%%%%%
\begin{document}

%%%%%%%%%%%%%%%%%%%%%%%%%%%%%%%%%% PORTADA %%%%%%%%%%%%%%%%%%%%%%%%%%%%%%%%%%%%%%%%%%%%
																					%%%
\begin{center}																		%%%
\newcommand{\HRule}{\rule{\linewidth}{0.5mm}}									%%%\left
 
 %%%

\begin{minipage}{0.48\textwidth} \begin{flushleft}
\includegraphics[scale = 0.30]{log_Eduim.jpg}
\end{flushleft}\end{minipage}
\begin{minipage}{0.48\textwidth} \begin{flushright}
\includegraphics[scale = 0.30]{log_uta.png}
\end{flushright}\end{minipage}

													 								%%%
\vspace*{1.0cm}								%%%
																					%%%	
\textsc{\huge Universidad de Tarapaca \vspace{5px}}\\[1.5cm]	

\textsc{\LARGE Escuela Universitaria de \\ \vspace{5px} Ingenieria Mecanica }\\[1.5cm]													%%%

\begin{minipage}{0.9\textwidth} 
\begin{center}																					%%%
\textsc{\LARGE Informe 02 }
\end{center}
\end{minipage}\\[0.5cm]
%%%
    																				%%%
 			\vspace*{1cm}																		%%%
																					%%%
\HRule \\[0.4cm]																	%%%
{ \huge \bfseries Mini Impresora 3D, usando servo motores}\\[0.4cm]	%%%
 																					%%%
\HRule \\[1.5cm]																	%%%
 																				%%%
																					%%%
\begin{minipage}{0.46\textwidth}													%%%
\begin{flushleft} \large															%%%

% Aqui a continuación pongan los nombres de los integrantes
\emph{Estudiantes:}\\	
Cristian Espinoza J.\\
Francisco Saez\\
Brain Pizarro
%%%
			%\vspace*{2cm}	
            													%%%
										 						%%%
\end{flushleft}																		%%%
\end{minipage}		
																%%%
\begin{minipage}{0.52\textwidth}		
\vspace{-0.6cm}											%%%
\begin{flushright} \large															%%%
\emph{Profesor:} \\																	%%%
Dr. Cristobal Castro\\														%%%
\end{flushright}																	%%%
\end{minipage}	
\vspace*{1cm}
%\begin{flushleft}
 	
%\end{flushleft}
%%%
 		\flushleft{\textbf{\Large Intro. a la Ingeniería}\\II001}\\																		%%%
\vspace{2cm} 
\end{center}

\begin{center}	
Arica,\\
{\large \today}																	%%%
\end{center}										  						
%%%%%%%%%%%%%%%% FINAL DE LA PORTADA %%%%%%%%%%%%%%%%%%%%%%%%%%%


\tableofcontents 

\newpage
%%%%%%%%%%%%%%%%%%%%%%
\section{Resumen.}



Nuestro proyecto es crear un robot armable que sirva como un medio para introducir, a los niños de 6 años y  más, a la robótica. A través de este robot buscamos enseñar a los niños a armar piezas mecánicas, circuitos eléctricos y a programar sus propios robots. Nuestro equipo esta formado enteramente por estudiantes de ingeniería mecánica y mecatrónica y estamos todos cursando el primer semestre. esto es importante para nosotros porque lo que hemos hecho en este proyecto no se compara con las cosas que hemos aprendido.


%%%%%%%%%%%%%%%%%%%%%%%
\section{Introducción.} 

Solo hace veinte años atrás los niños jugaban afuera de sus casas, andaban en bicicletas, construían fortificaciones, jugaban al football o inventaban juegos jugando con la imaginación. Sin duda que los niños de antes no requerían de dispositivos electrónicos o el uso de la Internet para desarrollar la imaginación.

Las familias de hoy son diferentes y la tecnología del siglo 21 ha moldeado en gran medida el que hacer en nuestros hogares e incluso en las escuelas. Pero el estudio del impacto de la tecnología en los niños no es objetivo de este trabajo sino de invitar a los niños a usar su imaginación en la creación de juegos usando la tecnología y las herramientas disponibles en Internet.

Todos los días vemos a niños jugando con sus celulares y viendo videos en youtube. Pero cuantos de ellos usan sus propios juegos o usan la tecnología para crear sus propias aplicaciones. No lo sabemos.

Si usamos esa tecnología y le mostramos a los niños como usarla para crear sus propios juegos y aplicaciones o para crear maquinas o robots, sin duda que el celular y la Internet se convierten en un medio de investigación y desarrollo. 

En el proceso, como alumnos de ingeniería, también aprenderemos a desarrollar nuestras propias habilidades, adquirir conocimiento y fortalecer nuestras expectativas frente a la elección de estudiar ingeniería.

\vspace{1cm} 
\begin{figure}[H]
	\begin{center}
 		\includegraphics[width = 0.5\textwidth]{robot1.jpg}
 		\captionof{figure}{\label{fig:robot I2}Modelación del Robot I2 en TinkerCad} 
	\end{center} 
\end{figure}

\vspace{2cm} 

%%%%%%%%%%%%%%%%%%%%%
\section{Identificación del Problema.}


Se requiere diseñar un proyecto que debe ser atractivo y no muy complejo para ser desarrollado por niños de 6 años de edad y más. La gran mayoría de niños son usuarios de la tecnología. Ellos tienen acceso casi instantáneo a plataformas como Youtube, Instagram, juegos online, etc. comparten fotos, mensajes de textos y videos. Pero estamos seguro que todos los niños tienen el potencial para desarrollar sus propias aplicaciones o para usar su expertice tecnológico para crear robots y convertirse en creadores.

\begin{itemize}
\item  El proyecto debe desarrollar habilidades motoras y cognitivas. Por ejemplo; puede armar elementos mecánicos y desarrollar un programa para entender el pensamiento lógico de programación.
\item  La tecnología a usar debe ser reconocida por los niños. Se debe incorporar herramientas como: el celular, computadores, iPads, herramientas que ya son de uso frecuente por los niños. 
\item  La duración del proyecto no debe ser extenso, puesto que los niños adquieren información en forma casi automática. Y se aburren con tareas muy largas o pierden el interés.
\item El proyecto debe ser atractivo visualmente, para llamar la atención de los niños. 
\item En las etapas del proyecto, éste debe introducir nuevas tecnologías, por ejemplo; programar Arduino, fuentes de energía, crear aplicaciones en Android, etc.

\end{itemize}


%%%%%%%%%%%%%%%%%%%%%%%%
\section{Posibles Soluciones.}

A continuación se dan a conocer posibles proyectos que puedan dar solución a la falta de interés o desconocimientos de nuevas tecnologías para desarrollar la creatividad y así los niños puedan pasar del modo usuario al modo creadores.

\begin{itemize}
\item  Llevar a los niños a exposiciones de tecnología, por ejemplo; ferias de robótica y de ciencias.
\item  Entregarles a los niños un kit de robótica para que ellos lo puedan armar. Piezas mecánicas y eléctricas. Y también lo puedan programar.
\item  Creación de academias de programación y robótica en las escuelas.
\item  Competencias de proyectos de innovación, ciencia y tecnología a nivel escolar. 
\item  Hacer campamentos de verano de robótica y programación.

\end{itemize}


\begin{figure}[H]
	\begin{center}
 		\includegraphics[width = 0.5\textwidth]{Robot2.jpg}
 		\captionof{figure}{\label{fig:Camprobotics}Campamento de robótica en Asia (foto de campasia.asia)} 
	\end{center} 
\end{figure}


%%%%%%%%%%%%%%%%%%%%%%%%%%
\subsection{Solución.}

Creemos que la mejor solución, no la única, es desarrollar un proyecto de Kit de robótica donde puedan desarrollar habilidades mecánicas, eléctricas y de programación. 

Nuestro objetivo es crear un robot que sea armable, que se conecte a circuitos eléctricos a través de un controlador y que sea fácil de programar para que los niños sientan el desafío de crear algo por ellos mismos. 
%%
\vspace{0.5cm} 
\\Algunas de las ventajas que adquieren los niños involucrados en este proyecto en comparación con aquellos que sólo son usuarios, son enumeradas abajo:

\begin{enumerate}
\item  Reconocer y armar piezas mecánicas,
\item  Reconocer y armar piezas eléctricas,
\item  Conexión de cables eléctricos por medio de un diagrama de conexión,
\item  Reconocer fuentes de energías (baterias),
\item  Conocer controladores como Arduino Uno,
\item  Programar Arduino Uno,
\item  Usar programación en bloques para Arduino,
\item  Reconocer tecnología de Bluetooth y 
\item  Conexiones por Bluetooth.
\end{enumerate}

\vspace{0.5cm} 
Existen algunas desventajas con respecto al diseño que deben ser consideradas, estas son:

\begin{enumerate}
\item Para una mayor autonomía, el robot debe usar una buena batería y debe ser recargable.
\item El diseño debe ser atractivo, para provocar interés en el niño.
\item La programación debe ser por bloques, como por ejemplo: Scratch.
\item El instructivo de armado debe ser ilustrativo, junto con la configuración del cableado.
\end{enumerate}

\subsection{Prototipo Feo}

El prototipo esta formado por los siguientes elementos:

\begin{enumerate}
\item 01 Arduino Uno R3
\item 01 Protoboard
\item 01 Puente H
\item 02 Motores DC
\item Cables Jumpers
\item 01 Rueda loca
\item Batería recargable
\item Bluetooth
\end{enumerate}

\begin{figure}[H]
	\begin{center}
 		\includegraphics[width = 0.4\textwidth]{RobotI2.jpg}
 		\captionof{figure}{\label{fig:RobotI2}Primer prototipo del I2,}
   \vspace{1cm} 
        \includegraphics[width = 0.6\textwidth]{circuito.png}
        \captionof{figure}{\label{fig:circuito}Circuito eléctrico de I2}
	\end{center} 
\end{figure}

\section{Conclusiones.}

Esta experiencia ha sido la más gratificante hasta ahora en la universidad. Realmente pudimos introducirnos a la ingeniería con conceptos de diseño y procesos de fabricación aunque sean a muy baja escala. aprendimos a descubrir el planteamiento de un problema y ha establecer las posibles soluciones. 

En el proceso algunos aprendimos una pincelada (introducción) a programas 3D, programación básica en Arduino, circuitos eléctricos, corte con máquina laser, a medir con pie de metro, etc.

  \vspace{1cm}
\begin{figure}[H]
	\begin{center}
 		\includegraphics[width = 0.5\textwidth]{photo1.jpg}
 		\captionof{figure}{\label{fig:robot I2}Instalación de circuitos en puente H} 
	\end{center} 
\end{figure}
   \vspace{1cm} 
\begin{figure}[H]
	\begin{center}
 		\includegraphics[width = 0.3\textwidth]{photo3.jpg}
 		\captionof{figure}{\label{fig:robot I2}Subiendo el programa al Arduino} 
	\end{center} 
\end{figure}

\vspace{1.5cm}
Con respecto al diseño del proyecto, cumplió con todas nuestras expectativas. Logramos fabricar el prototipo feo que habíamos proyectado en CAD. Sin duda es un prototipo que requiere de mejoras, tanto en el diseño estético como en el funcional. Todavía el I2 (nuestro robot) esta en etapa de prueba y requiere mejorar sus características para que sea atractivo para los niños.


\begin{enumerate}
\item Agregar un sensor de Ultra sonido
\item Lograr usar un bluetooth o wifi para manejar el carro
\item Mejorar el diseño estético del carro.
\item Hacer que el carro sea más liviano, entre otras cosas.

\end{enumerate}







\end{document}

